% =========================================
% PHẦN 2: Quá trình Quyết định Markov (MDP)
% =========================================

\begin{frame}
  \centering
  \Huge \textbf{PHẦN 2} \\[0.5cm]
  \LARGE Quá trình Quyết định Markov (MDP)
\end{frame}

%---------------------------------
\begin{frame}{Quá trình Quyết định Markov (MDP) (1)}
\begin{itemize}
    \item \textbf{Định nghĩa:} MDP là khung toán học cho việc ra quyết định trong môi trường ngẫu nhiên.
    \item \textbf{Thành phần của MDP:}
    \begin{itemize}
        \item $S$: Không gian trạng thái
        \item $A$: Không gian hành động
        \item $P(s'|s,a)$: Hàm xác suất chuyển trạng thái — xác suất chuyển từ trạng thái $s$ sang $s'$ sau khi thực hiện hành động $a$.
        \item $R(s,a)$: Hàm phần thưởng — đánh giá mức độ hiệu quả của hành động $a$ tại trạng thái $s$.
    \end{itemize}
    \item \textbf{Tính chất Markov:}
    \[
    P(s_{t+1} | s_t, a_t) = P(s_{t+1} | s_0, a_0, \ldots, s_t, a_t)
    \]
    Trạng thái tương lai \textbf{chỉ phụ thuộc} vào trạng thái hiện tại, không phụ thuộc lịch sử trước đó.
\end{itemize}
\end{frame}

%---------------------------------
\begin{frame}{Quá trình Quyết định Markov (MDP) (2)}
\begin{itemize}
    \item \textbf{Mục tiêu:} Tìm ra chính sách tối ưu $\pi^*(a|s)$ ánh xạ trạng thái sang hành động nhằm tối đa hóa phần thưởng chiết khấu lũy kế kỳ vọng:
    \[
    \pi^* = \arg\max_\pi\; \mathbb{E}\left[\sum_{t=0}^{\infty} \gamma^t R(s_t, a_t)\right]
    \]
    \item Trong đó $\gamma \in (0,1]$ là \textbf{hệ số chiết khấu} — đại diện cho tầm quan trọng tương đối của phần thưởng tương lai so với phần thưởng tức thì.
    \item Dự án sử dụng $\gamma = 0.99$ — ưu tiên tối ưu hóa dài hạn trong suốt 1 giờ mô phỏng.
\end{itemize}
\end{frame}

%---------------------------------
\begin{frame}{Ánh xạ bài toán đèn giao thông sang MDP (1)}
\begin{itemize}
    \item \textbf{Tác nhân (Agent):} Hệ thống điều khiển tín hiệu giao thông (DQN agent)
    \item \textbf{Môi trường:} Ngã tư 4 hướng được mô phỏng bởi \textbf{SUMO}, giao tiếp qua \textbf{TraCI API}
    \item \textbf{Trạng thái ($s$):} Vector liên tục thu thập từ 8 làn xe vào nút
    \begin{itemize}
        \item Mỗi làn được biểu diễn bởi \textbf{4 đặc trưng}:
            \begin{itemize}
                \item Số xe dừng — queue length ($q_i$)
                \item Tốc độ trung bình — average speed ($v_i$)
                \item Thời gian chờ tích lũy — waiting time ($w_i$)
                \item Số xe quy đổi PCU — weighted vehicle count ($c_i$)
            \end{itemize}
        \item 4 hướng $\times$ 2 làn/hướng $\times$ 4 đặc trưng $= s \in \mathbb{R}^{32}$
        \item \textbf{Hệ số PCU Việt Nam:} xe máy (0.5), ô tô (1.5), xe buýt/tải (2.0)
        \item $q_i$, $w_i$, $c_i$ được nhân thêm \textbf{direction weight} theo hướng xe để agent ưu tiên hướng có lưu lượng lớn hơn
    \end{itemize}
\end{itemize}
\end{frame}

%---------------------------------
\begin{frame}{Ánh xạ bài toán đèn giao thông sang MDP (2)}
\begin{itemize}
    \item \textbf{Không gian hành động ($A$):} 4 pha tín hiệu được định nghĩa trước
    \begin{itemize}
        \item $a_0$ (Pha 0): Bắc -- Nam \textbf{xanh} — cho phép xe từ hướng Bắc/Nam
        \item $a_1$ (Pha 1): Bắc -- Nam \textbf{vàng} — pha giải phóng trước khi chuyển sang Đông/Tây
        \item $a_2$ (Pha 2): Đông -- Tây \textbf{xanh} — cho phép xe từ hướng Đông/Tây
        \item $a_3$ (Pha 3): Đông -- Tây \textbf{vàng} — pha giải phóng trước khi chuyển sang Bắc/Nam
    \end{itemize}
    \item Agent ra quyết định mỗi \textbf{5 giây} mô phỏng (\texttt{action\_duration = 5s})
    \item Kịch bản lưu lượng bất đối xứng: \textbf{Đông--Tây : Bắc--Nam = 2:1}\\
          EW xanh tối thiểu giữ 60s; NS xanh tối thiểu giữ 30s
\end{itemize}
\end{frame}

%---------------------------------
\begin{frame}{Ràng buộc Thời gian Pha (Per-Phase Constraint)}
\begin{itemize}
    \item \textbf{Vấn đề:} Nếu agent chuyển pha tùy ý, có thể gây hỗn loạn:
    \begin{itemize}
        \item Pha xanh chỉ kéo 5 giây → phương tiện vừa vào giao lộ không kịp qua
        \item Không ổn định trong quá trình học → reward dao động lớn
    \end{itemize}

    \item \textbf{Giải pháp:} Ràng buộc riêng biệt cho từng pha xanh:
    \begin{itemize}
        \item \textbf{Pha 2 — EW xanh (trục chính):} $t_{\min} = 60$s, $t_{\max} = 140$s
        \item \textbf{Pha 0 — NS xanh (trục phụ):} $t_{\min} = 30$s, $t_{\max} = 70$s
    \end{itemize}

    \item \textbf{Cơ chế thi hành:}
    \begin{itemize}
        \item Chuyển pha trước $t_{\min}$: \textbf{bác bỏ hành động} + phạt reward $-1.0$
        \item Giữ pha vượt $t_{\max}$: \textbf{bắt buộc chuyển} sang pha kế tiếp
    \end{itemize}

    \item \textbf{Hiệu ứng bất đối xứng có chủ ý:} EW có $t_{\min}$ gấp đôi NS → agent tự nhiên học ưu tiên trục chính, phù hợp tỉ lệ lưu lượng 2:1.
\end{itemize}
\end{frame}

%---------------------------------
\begin{frame}{Ánh xạ bài toán đèn giao thông sang MDP (3)}
\begin{itemize}
    \item \textbf{Hàm phần thưởng ($R$):}
    \begin{itemize}
        \item Phần thưởng được tính dựa trên mức phạt từ tình trạng ùn tắc hiện tại:
        \[
        R_t = -\left( Q_t + \frac{W_t}{60} \right) + \text{penalty}
        \]
        \begin{itemize}
            \item $Q_t$: tổng số xe dừng (queue) trên tất cả làn, có nhân direction weight
            \item $W_t$: tổng thời gian chờ trên tất cả làn, chia 60 để chuẩn hóa đơn vị
            \item $\text{penalty} = -1.0$: áp dụng khi agent cố chuyển pha trước $t_{\min}$
        \end{itemize}
        \item $R_t$ càng gần 0: giao thông càng thông thoáng
        \item $R_t$ càng âm: ùn tắc càng nghiêm trọng
    \end{itemize}

    \item \textbf{Chuyển đổi trạng thái ($P$):}
    \begin{itemize}
        \item Phụ thuộc vào hành động $a_t$ (pha đèn được áp dụng)
        \item Động lực học phương tiện trong mô phỏng SUMO
        \item Sự xuất hiện ngẫu nhiên của xe mới theo phân bố xác suất cố định
    \end{itemize}
\end{itemize}
\end{frame}
